\IfFileExists{stacks-project.cls}{%
\documentclass{stacks-project}
}{%
\documentclass{amsart}
}

% For dealing with references we use the comment environment
\usepackage{verbatim}
\newenvironment{reference}{\comment}{\endcomment}
%\newenvironment{reference}{}{}
\newenvironment{slogan}{\comment}{\endcomment}
\newenvironment{history}{\comment}{\endcomment}

% For commutative diagrams we use Xy-pic
\usepackage[all]{xy}

% We use 2cell for 2-commutative diagrams.
\xyoption{2cell}
\UseAllTwocells

% We use multicol for the list of chapters between chapters
\usepackage{multicol}

% This is generally recommended for better output
\usepackage{lmodern}
\usepackage[T1]{fontenc}

% For cross-file-references
\usepackage{xr-hyper}

% Package for hypertext links:
\usepackage{hyperref}

% For any local file, say "hello.tex" you want to link to please
% use \externaldocument[hello-]{hello}
\externaldocument[introduction-]{introduction}
\externaldocument[conventions-]{conventions}
\externaldocument[sets-]{sets}
\externaldocument[categories-]{categories}
\externaldocument[topology-]{topology}
\externaldocument[sheaves-]{sheaves}
\externaldocument[sites-]{sites}
\externaldocument[stacks-]{stacks}
\externaldocument[fields-]{fields}
\externaldocument[algebra-]{algebra}
\externaldocument[brauer-]{brauer}
\externaldocument[homology-]{homology}
\externaldocument[derived-]{derived}
\externaldocument[simplicial-]{simplicial}
\externaldocument[more-algebra-]{more-algebra}
\externaldocument[smoothing-]{smoothing}
\externaldocument[modules-]{modules}
\externaldocument[sites-modules-]{sites-modules}
\externaldocument[injectives-]{injectives}
\externaldocument[cohomology-]{cohomology}
\externaldocument[sites-cohomology-]{sites-cohomology}
\externaldocument[dga-]{dga}
\externaldocument[dpa-]{dpa}
\externaldocument[sdga-]{sdga}
\externaldocument[hypercovering-]{hypercovering}
\externaldocument[schemes-]{schemes}
\externaldocument[constructions-]{constructions}
\externaldocument[properties-]{properties}
\externaldocument[morphisms-]{morphisms}
\externaldocument[coherent-]{coherent}
\externaldocument[divisors-]{divisors}
\externaldocument[limits-]{limits}
\externaldocument[varieties-]{varieties}
\externaldocument[topologies-]{topologies}
\externaldocument[descent-]{descent}
\externaldocument[perfect-]{perfect}
\externaldocument[more-morphisms-]{more-morphisms}
\externaldocument[flat-]{flat}
\externaldocument[groupoids-]{groupoids}
\externaldocument[more-groupoids-]{more-groupoids}
\externaldocument[etale-]{etale}
\externaldocument[chow-]{chow}
\externaldocument[intersection-]{intersection}
\externaldocument[pic-]{pic}
\externaldocument[weil-]{weil}
\externaldocument[adequate-]{adequate}
\externaldocument[dualizing-]{dualizing}
\externaldocument[duality-]{duality}
\externaldocument[discriminant-]{discriminant}
\externaldocument[derham-]{derham}
\externaldocument[local-cohomology-]{local-cohomology}
\externaldocument[algebraization-]{algebraization}
\externaldocument[curves-]{curves}
\externaldocument[resolve-]{resolve}
\externaldocument[models-]{models}
\externaldocument[functors-]{functors}
\externaldocument[equiv-]{equiv}
\externaldocument[pione-]{pione}
\externaldocument[etale-cohomology-]{etale-cohomology}
\externaldocument[proetale-]{proetale}
\externaldocument[relative-cycles-]{relative-cycles}
\externaldocument[more-etale-]{more-etale}
\externaldocument[trace-]{trace}
\externaldocument[crystalline-]{crystalline}
\externaldocument[spaces-]{spaces}
\externaldocument[spaces-properties-]{spaces-properties}
\externaldocument[spaces-morphisms-]{spaces-morphisms}
\externaldocument[decent-spaces-]{decent-spaces}
\externaldocument[spaces-cohomology-]{spaces-cohomology}
\externaldocument[spaces-limits-]{spaces-limits}
\externaldocument[spaces-divisors-]{spaces-divisors}
\externaldocument[spaces-over-fields-]{spaces-over-fields}
\externaldocument[spaces-topologies-]{spaces-topologies}
\externaldocument[spaces-descent-]{spaces-descent}
\externaldocument[spaces-perfect-]{spaces-perfect}
\externaldocument[spaces-more-morphisms-]{spaces-more-morphisms}
\externaldocument[spaces-flat-]{spaces-flat}
\externaldocument[spaces-groupoids-]{spaces-groupoids}
\externaldocument[spaces-more-groupoids-]{spaces-more-groupoids}
\externaldocument[bootstrap-]{bootstrap}
\externaldocument[spaces-pushouts-]{spaces-pushouts}
\externaldocument[spaces-chow-]{spaces-chow}
\externaldocument[groupoids-quotients-]{groupoids-quotients}
\externaldocument[spaces-more-cohomology-]{spaces-more-cohomology}
\externaldocument[spaces-simplicial-]{spaces-simplicial}
\externaldocument[spaces-duality-]{spaces-duality}
\externaldocument[formal-spaces-]{formal-spaces}
\externaldocument[restricted-]{restricted}
\externaldocument[spaces-resolve-]{spaces-resolve}
\externaldocument[formal-defos-]{formal-defos}
\externaldocument[defos-]{defos}
\externaldocument[cotangent-]{cotangent}
\externaldocument[examples-defos-]{examples-defos}
\externaldocument[algebraic-]{algebraic}
\externaldocument[examples-stacks-]{examples-stacks}
\externaldocument[stacks-sheaves-]{stacks-sheaves}
\externaldocument[criteria-]{criteria}
\externaldocument[artin-]{artin}
\externaldocument[quot-]{quot}
\externaldocument[stacks-properties-]{stacks-properties}
\externaldocument[stacks-morphisms-]{stacks-morphisms}
\externaldocument[stacks-limits-]{stacks-limits}
\externaldocument[stacks-cohomology-]{stacks-cohomology}
\externaldocument[stacks-perfect-]{stacks-perfect}
\externaldocument[stacks-introduction-]{stacks-introduction}
\externaldocument[stacks-more-morphisms-]{stacks-more-morphisms}
\externaldocument[stacks-geometry-]{stacks-geometry}
\externaldocument[moduli-]{moduli}
\externaldocument[moduli-curves-]{moduli-curves}
\externaldocument[examples-]{examples}
\externaldocument[exercises-]{exercises}
\externaldocument[guide-]{guide}
\externaldocument[desirables-]{desirables}
\externaldocument[coding-]{coding}
\externaldocument[obsolete-]{obsolete}
\externaldocument[fdl-]{fdl}
\externaldocument[index-]{index}

% Theorem environments.
%
\theoremstyle{plain}
\newtheorem{theorem}[subsection]{Theorem}
\newtheorem{proposition}[subsection]{Proposition}
\newtheorem{lemma}[subsection]{Lemma}

\theoremstyle{definition}
\newtheorem{definition}[subsection]{Definition}
\newtheorem{example}[subsection]{Example}
\newtheorem{exercise}[subsection]{Exercise}
\newtheorem{situation}[subsection]{Situation}

\theoremstyle{remark}
\newtheorem{remark}[subsection]{Remark}
\newtheorem{remarks}[subsection]{Remarks}

\numberwithin{equation}{subsection}

% Macros
%
\def\lim{\mathop{\mathrm{lim}}\nolimits}
\def\colim{\mathop{\mathrm{colim}}\nolimits}
\def\Spec{\mathop{\mathrm{Spec}}}
\def\Hom{\mathop{\mathrm{Hom}}\nolimits}
\def\Ext{\mathop{\mathrm{Ext}}\nolimits}
\def\SheafHom{\mathop{\mathcal{H}\!\mathit{om}}\nolimits}
\def\SheafExt{\mathop{\mathcal{E}\!\mathit{xt}}\nolimits}
\def\Sch{\mathit{Sch}}
\def\Mor{\mathop{\mathrm{Mor}}\nolimits}
\def\Ob{\mathop{\mathrm{Ob}}\nolimits}
\def\Sh{\mathop{\mathit{Sh}}\nolimits}
\def\NL{\mathop{N\!L}\nolimits}
\def\CH{\mathop{\mathrm{CH}}\nolimits}
\def\proetale{{pro\text{-}\acute{e}tale}}
\def\etale{{\acute{e}tale}}
\def\QCoh{\mathit{QCoh}}
\def\Ker{\mathop{\mathrm{Ker}}}
\def\Im{\mathop{\mathrm{Im}}}
\def\Coker{\mathop{\mathrm{Coker}}}
\def\Coim{\mathop{\mathrm{Coim}}}

% Boxtimes
%
\DeclareMathSymbol{\boxtimes}{\mathbin}{AMSa}{"02}

%
% Macros for moduli stacks/spaces
%
\def\QCohstack{\mathcal{QC}\!\mathit{oh}}
\def\Cohstack{\mathcal{C}\!\mathit{oh}}
\def\Spacesstack{\mathcal{S}\!\mathit{paces}}
\def\Quotfunctor{\mathrm{Quot}}
\def\Hilbfunctor{\mathrm{Hilb}}
\def\Curvesstack{\mathcal{C}\!\mathit{urves}}
\def\Polarizedstack{\mathcal{P}\!\mathit{olarized}}
\def\Complexesstack{\mathcal{C}\!\mathit{omplexes}}
% \Pic is the operator that assigns to X its picard group, usage \Pic(X)
% \Picardstack_{X/B} denotes the Picard stack of X over B
% \Picardfunctor_{X/B} denotes the Picard functor of X over B
\def\Pic{\mathop{\mathrm{Pic}}\nolimits}
\def\Picardstack{\mathcal{P}\!\mathit{ic}}
\def\Picardfunctor{\mathrm{Pic}}
\def\Deformationcategory{\mathcal{D}\!\mathit{ef}}
\usepackage{cleveref}
\crefname{originallemma}{original lemma}{original lemmas}
\Crefname{originallemma}{Original lemma}{Original lemmas}
\makeatletter
\newlabel{lemma-bounded-fibres}{{03JX}{source}{}{sourcepointer-1}{}}
\newlabel{lemma-bounded-fibres@cref}{{[originallemma][0][]03JX}{[1][0][]source}{}{}{}}
\makeatother
\crefname{originallemma}{original lemma}{original lemmas}
\Crefname{originallemma}{Original lemma}{Original lemmas}
\makeatletter
\newlabel{lemma-extend-integral-morphism}{{0G2D}{source}{}{sourcepointer-2}{}}
\newlabel{lemma-extend-integral-morphism@cref}{{[originallemma][0][]0G2D}{[1][0][]source}{}{}{}}
\makeatother
\crefname{originallemma}{original lemma}{original lemmas}
\Crefname{originallemma}{Original lemma}{Original lemmas}
\makeatletter
\newlabel{lemma-quasi-compact-reasonable-stratification}{{07S8}{source}{}{sourcepointer-3}{}}
\newlabel{lemma-quasi-compact-reasonable-stratification@cref}{{[originallemma][0][]07S8}{[1][0][]source}{}{}{}}
\makeatother
\crefname{originalreference}{original reference}{original references}
\Crefname{originalreference}{Original reference}{Original references}
\makeatletter
\newlabel{spaces-morphisms-lemma-quasi-finite-separated-quasi-affine}{{0ABS}{source}{}{sourcepointer-4}{}}
\newlabel{spaces-morphisms-lemma-quasi-finite-separated-quasi-affine@cref}{{[originalreference][0][]0ABS}{[1][0][]source}{}{}{}}
\makeatother
\crefname{originalreference}{original reference}{original references}
\Crefname{originalreference}{Original reference}{Original references}
\makeatletter
\newlabel{spaces-morphisms-lemma-semi-diagonal}{{03KO}{source}{}{sourcepointer-5}{}}
\newlabel{spaces-morphisms-lemma-semi-diagonal@cref}{{[originalreference][0][]03KO}{[1][0][]source}{}{}{}}
\makeatother
\title{A quasi-compact quasi-separated algebraic space admits a surjective integral cover by a scheme factoring locally through a given etale cover}
\author{The Stacks Project Authors}
\date{stacks-project:decent-spaces.tex:local-snapshot-20261012; GNU FDL 1.2 or later}
\begin{document}
\maketitle
\noindent\textit{Editorial note.} No omission markers in-span.\par\medskip
\noindent\textit{Allowed standard prerequisites (editorial).}\begin{itemize}
\item Bounded fibres of the etale cover. lemma-bounded-fibres (carded L26703). The degree bound d for U -> X.
\item Degree stratification. lemma-quasi-compact-reasonable-stratification (carded L26710). The degree of the fibres of a quasi-compact etale morphism does not go up under specialization.
\item Integral extension of the cover domain. lemma-extend-integral-morphism (same chapter, carded L26059). Factorization Z -> Z' -> Y with Z' -> Y integral.
\item Normalization in a quasi-finite separated morphism. Morphisms of Spaces, lemma quasi-finite-separated-quasi-affine (cross-chapter, named). The factorization U -> Y -> X with pi integral.
\item Scheme-theoretic density and integral composition. Morphisms of Spaces, lemmas semi-diagonal and composition-integral (cross-chapter, named). The decomposition U x\_X Y = U amal W and integrality of the assembled morphism.
\end{itemize}\medskip
\ifdefined\hypertarget\hypertarget{sourcepointer-1}{}\fi
\noindent\textit{Source reference.} Original lemma 03JX: sources/stacks-decent-spaces/pinned-original-source.tex, source lines 952--952. Navigation pointer only: the counted proof cites this statement by \textbackslash{}ref as a standing input; the deduction the counted proof has to make is carried out inside the counted span, and no part of that deduction is delegated to this pointer.\par\medskip
\ifdefined\hypertarget\hypertarget{sourcepointer-2}{}\fi
\noindent\textit{Source reference.} Original lemma 0G2D: sources/stacks-decent-spaces/pinned-original-source.tex, source lines 2002--2002. Navigation pointer only: the counted proof cites this statement by \textbackslash{}ref as a standing input; the deduction the counted proof has to make is carried out inside the counted span, and no part of that deduction is delegated to this pointer.\par\medskip
\ifdefined\hypertarget\hypertarget{sourcepointer-3}{}\fi
\noindent\textit{Source reference.} Original lemma 07S8: sources/stacks-decent-spaces/pinned-original-source.tex, source lines 1653--1653. Navigation pointer only: the counted proof cites this statement by \textbackslash{}ref as a standing input; the deduction the counted proof has to make is carried out inside the counted span, and no part of that deduction is delegated to this pointer.\par\medskip
\ifdefined\hypertarget\hypertarget{sourcepointer-4}{}\fi
\noindent\textit{Source reference.} Original reference 0ABS: focused-seeds/geometry/S9-decent-there-is-a-scheme-integral-over/ref-tags.txt, source lines 5--5. Navigation pointer only: the counted proof cites this statement by \textbackslash{}ref as a standing input; the deduction the counted proof has to make is carried out inside the counted span, and no part of that deduction is delegated to this pointer.\par\medskip
\ifdefined\hypertarget\hypertarget{sourcepointer-5}{}\fi
\noindent\textit{Source reference.} Original reference 03KO: focused-seeds/geometry/S9-decent-there-is-a-scheme-integral-over/ref-tags.txt, source lines 6--6. Navigation pointer only: the counted proof cites this statement by \textbackslash{}ref as a standing input; the deduction the counted proof has to make is carried out inside the counted span, and no part of that deduction is delegated to this pointer.\par\medskip
\begin{lemma}
\label{lemma-there-is-a-scheme-integral-over}
Let $S$ be a scheme. Let $X$ be a quasi-compact and quasi-separated
algebraic space over $S$.
\begin{enumerate}
\item There exists a surjective integral morphism $Y \to X$ where $Y$
is a scheme,
\item given a surjective \'etale morphism $U \to X$ we may choose
$Y \to X$ such that for every $y \in Y$ there is an open neighbourhood
$V \subset Y$ such that $V \to X$ factors through $U$.
\end{enumerate}
\end{lemma}

\begin{proof}
Part (1) is the special case of part (2) where $U = X$.
Choose a surjective \'etale morphism $U' \to U$
where $U'$ is a scheme. It is clear that we may replace $U$ by $U'$
and hence we may assume $U$ is a scheme. Since $X$ is quasi-compact,
there exist finitely many affine opens $U_i \subset U$ such that
$U' = \coprod U_i \to X$ is surjective.
After replacing $U$ by $U'$ again, we see that we may assume $U$ is affine.
Since $X$ is quasi-separated, hence reasonable, there exists an integer
$d$ bounding the degree of the geometric fibres of $U \to X$
(see Lemma \ref{lemma-bounded-fibres}).
We will prove the lemma by induction on $d$ for all quasi-compact
and separated schemes $U$ mapping surjective and \'etale onto $X$.
If $d = 1$, then $U = X$ and the result holds with $Y = U$.
Assume $d > 1$.

\medskip\noindent
We apply Morphisms of Spaces, Lemma
\ref{spaces-morphisms-lemma-quasi-finite-separated-quasi-affine}
and we obtain a factorization
$$
\xymatrix{
U \ar[rr]_j \ar[rd] & & Y \ar[ld]^\pi \\
& X
}
$$
with $\pi$ integral and $j$ a quasi-compact open immersion. We may and do
assume that $j(U)$ is scheme theoretically dense in $Y$. Then $U \times_X Y$
is a quasi-compact, separated scheme (being integral over $U$) and we have
$$
U \times_X Y = U \amalg W
$$
Here the first summand is the image of $U \to U \times_X Y$
(which is closed by
Morphisms of Spaces, Lemma \ref{spaces-morphisms-lemma-semi-diagonal}
and open because it is \'etale as a morphism between
algebraic spaces \'etale over $Y$) and
the second summand is the (open and closed) complement.
The image $V \subset Y$ of $W$ is an open subspace containing
$Y \setminus U$.

\medskip\noindent
The \'etale morphism $W \to Y$ has geometric fibres of cardinality $< d$.
Namely, this is clear for geometric points of $U \subset Y$ by inspection.
Since $|U| \subset |Y|$ is dense, it holds for all geometric points of $Y$
by Lemma \ref{lemma-quasi-compact-reasonable-stratification}
(the degree of the fibres of a quasi-compact \'etale morphism
does not go up under specialization). Thus we may apply the induction
hypothesis to $W \to V$ and find a surjective integral morphism
$Z \to V$ with $Z$ a scheme, which Zariski locally factors through $W$.
Choose a factorization $Z \to Z' \to Y$ with $Z' \to Y$ integral and
$Z \to Z'$ open immersion
(Lemma \ref{lemma-extend-integral-morphism}).
After replacing $Z'$ by the scheme theoretic closure of $Z$ in $Z'$
we may assume that $Z$ is scheme theoretically dense in $Z'$.
After doing this we have $Z' \times_Y V = Z$. Finally,
let $T \subset Y$ be the induced closed subspace structure on $Y \setminus V$.
Consider the morphism
$$
Z' \amalg T \longrightarrow X
$$
This is a surjective integral morphism by construction.
Since $T \subset U$ it is clear that the morphism $T \to X$
factors through $U$. On the other hand, let $z \in Z'$
be a point. If $z \not \in Z$, then $z$ maps to a point of
$Y \setminus V \subset U$ and we find a neighbourhood of $z$
on which the morphism factors through $U$.
If $z \in Z$, then we have an open neighbourhood of $z$ in $Z$
(which is also an open neighbourhood of $z$ in $Z'$)
which factors through $W \subset U \times_X Y$ and hence through $U$.
\end{proof}

\end{document}
